% Part: sets-functions-relations
% Chapter: sets
% Section: russells-paradox

\documentclass[../../../include/open-logic-section]{subfiles}


\begin{document}

\olfileid{sfr}{set}{rus}
\olsection{రసెల్ వైరుధ్యం}

$\phi(x)$ అయ్యే $x$ల \emph{ఆ ఏకైక} సమితికి
$\Setabs{x}{\phi(x)}$ అనే సంకేతాన్ని వాడటానికి సమితుల సమానత్వ
సూత్రం ఆధారమిస్తుంది. అయితే ఈ సూత్రం \emph{నిజంగా}
అనుమతించేది కింది ఆలోచన మాత్రమే. $\phi$ ధర్మం గలవన్నీ,
అవి మాత్రమే సభ్యాలుగా గల సమితి \emph{ఉంటే}, \emph{అప్పుడు}
అటువంటి సమితి ఒక్కటే ఉంటుంది. మరో విధంగా చెప్పాలంటే,
ఏదైనా $\phi$ను నిర్ణయించాక $\Setabs{x}{\phi(x)}$ సమితి
\emph{ఉన్న పక్షంలో} అది ఏకైకమైనది.

కానీ ఈ షరతు చాలా ముఖ్యం!{} ప్రతి ధర్మం ద్వారా
\emph{ధర్మసంగ్రహం} సాధ్యం కాదు. అంటే కొన్ని ధర్మాలు సమితులను
నిర్వచించ\emph{వు}. ప్రతి ధర్మమూ సమితిని నిర్వచిస్తే,
స్పష్టమైన విరోధాలు ఎదురవుతాయి. ఇందుకు అత్యంత ప్రసిద్ధ ఉదాహరణ
రసెల్ వైరుధ్యం.

సమితులు ఇతర సమితులలో \tetoken{మూలకాలు}{!!{element}s} కావచ్చు. ఉదాహరణకు
$A$ అనే సమితి యొక్క ఘాత సమితి సమితులతోనే ఏర్పడుతుంది.
అందువల్ల ఒక సమితి మరొక సమితిలో \tetoken{ఒక మూలకం}{!!a{element}} అవుతుందా అని
ప్రశ్నించడం, పరిశీలించడం అర్థవంతమే. ఒక సమితి తనలో తానే సభ్యంగా
ఉండవచ్చా? సమితి అనే భావనలో దానిని నిషేధించేదేమీ కనిపించదు.
ఉదాహరణకు \emph{అన్ని} సమితులూ వస్తువుల సమూహంగా ఏర్పడితే,
వాటన్నింటినీ ఒకే సమితిలో చేర్చవచ్చని అనిపించవచ్చు: అదే అన్ని
సమితుల సమితి. అది కూడా ఒక సమితి కనుక, అన్ని సమితుల సమితిలో
\tetoken{ఒక మూలకం}{!!a{element}} అవుతుంది.

తనను తాను \tetoken{ఒక మూలకంగా}{!!a{element}} కలిగి ఉండకపోవడం, అంటే
\emph{తనలో తానే సభ్యంగా లేకపోవడం}, అనే ధర్మాన్ని పరిశీలించినప్పుడు
రసెల్ వైరుధ్యం వస్తుంది. తమను తాము \tetoken{ఒక మూలకంగా}{!!a{element}} కలిగి ఉండని
సమితులన్నింటి సమితి ఉందనుకుంటే ఏమవుతుంది?
\[
R = \Setabs{x}{x \notin x}
\]
అనే సమితి ఉందా? అది లేదని మనం నిరూపించగలం.

\begin{thm}[రసెల్ వైరుధ్యం]\ollabel{thm:russells-paradox}
	$R = \Setabs{x}{x \notin x}$ అనే సమితి లేదు.
\end{thm}

\begin{proof}
$R = \Setabs{x}{x \notin x}$ ఉంటే,
$R \in R$ కావడానికి అవసరమైన మరియు సరిపడే షరతు $R \notin R$
అవుతుంది. ఇది విరోధం.
\end{proof}

\begin{tagblock}{novice}
\begin{explain}
ఈ నిరూపణను నెమ్మదిగా పరిశీలిద్దాం. $R$ ఉంటే, $R \in R$ అవుతుందా
అని అడగవచ్చు. నిజంగానే $R \in R$ అనుకుందాం. తమలో తాము
\tetoken{మూలకాలు}{!!{element}s} కాని సమితులన్నింటి సమితిగా $R$ను నిర్వచించాం.
కనుక $R \in R$ అయితే, $R$ను నిర్వచించే ధర్మం $R$కే ఉండదు.
కానీ ఆ ధర్మం గల సమితులు మాత్రమే $R$లో ఉంటాయి. అందువల్ల
$R$ అనేది $R$లో \tetoken{ఒక మూలకం}{!!a{element}} కాలేదు; అంటే $R \notin R$.
అయితే $R$ అనేది $R$లో \tetoken{ఒక మూలకం}{!!a{element}} కావడమూ కాకపోవడమూ
రెండూ జరగలేవు. కాబట్టి విరోధం వచ్చింది.

$R \in R$ అనే ఊహ విరోధానికి దారితీసింది కనుక $R \notin R$.
కానీ ఇదీ విరోధానికే దారితీస్తుంది!{} ఎందుకంటే $R \notin R$ అయితే,
$R$ను నిర్వచించే ధర్మం $R$కే ఉంటుంది. కాబట్టి తమలో తాము సభ్యాలుగా
లేని ఇతర సమితులన్నింటిలాగే, $R$ కూడా $R$లో \tetoken{ఒక మూలకం}{!!a{element}} అవుతుంది.
మళ్లీ, అది $R$లో \tetoken{ఒక మూలకం}{!!a{element}} కాకపోవడమూ కావడమూ
రెండూ జరగలేవు.
\end{explain}
\end{tagblock}

\begin{digress}
రసెల్ వైరుధ్యంలో చిక్కుకోని సమితి సిద్ధాంతాన్ని ఎలా నిర్మించాలి?
అంటే $R = \Setabs{x}{x \notin x}$ ఉందనే \emph{అసంగతమైన}
వాదనకు దారితీయకుండా ఎలా ఉండాలి? సమితులు ఎప్పుడు ఉంటాయో,
ఎప్పుడు ఉండవో చెప్పడానికి అత్యంత కచ్చితమైన షరతులను ఇచ్చే
స్వీకృతాలను నిర్దేశించాలి.

ఈ అధ్యాయంలో వివరించిన సమితి సిద్ధాంతం అలా చేయదు. ఇది
\emph{నిజంగానే అనౌపచారికమైనది}. సమితులు సమితుల సమానత్వ సూత్రానికి
లోబడతాయని, కొన్ని సమితులు ఉన్నప్పుడు వాటి సమ్మేళనం, ఛేదనం మొదలైనవి
ఏర్పరచవచ్చని మాత్రమే ఇది చెబుతుంది. సమితి సిద్ధాంతాన్ని ఇంతకంటే
కచ్చితంగా అభివృద్ధి చేయవచ్చు.
\oliflabeldef{cumul:::part}{ఆ కచ్చితమైన విధానాన్ని
\olref[cumul][][]{part} భాగంలో పరిశీలిస్తాం. ప్రస్తుతానికి
అనౌపచారికంగానే ముందుకు సాగుతూ, విరోధాలను జాగ్రత్తగా తప్పించుకోవడానికి
ప్రయత్నిస్తాం.}{}
\end{digress}


\end{document}
